\documentclass[11pt,aspectratio=169]{beamer}
% Deliberately simple: default theme, no Metropolis.
\usetheme{default}
\usecolortheme{dove}
\usefonttheme{serif}
\setbeamertemplate{navigation symbols}{}
\setbeamertemplate{footline}[frame number]
\setbeamertemplate{itemize items}[circle]
\usepackage{amsmath,amssymb}
\usepackage{xcolor}
\definecolor{key}{rgb}{0.10,0.35,0.65}
\newcommand{\key}[1]{\textcolor{key}{\textbf{#1}}}
\newcommand{\bhat}{\hat{\mathbf{B}}}
\setbeamercolor{frametitle}{fg=key}
\newcommand{\qq}[1]{\begin{center}\emph{#1}\end{center}}

\title{Spin-Polarized Fusion Neutronics}
\subtitle{A from-scratch primer on everything we discussed\\{\small (assumes: multivariable calculus, linear algebra, intro upper-division physics)}}
\author{built for the train ride home}
\date{\today}

\begin{document}
\frame{\titlepage}

\begin{frame}{How to read this}
\begin{itemize}
\item Each part starts from something you already know and builds one step at a time.
\item \key{Blue} words are the terms worth remembering.
\item Formulas are here, but every one comes with a plain-English sentence. If you skip the
algebra, read the sentence.
\item Goal: by the end you can explain the whole project to a friend --- the physics, the
engineering, the computation, and why any of it matters.
\end{itemize}
\end{frame}

\begin{frame}{The one-sentence version}
\qq{If you line up the spins of fusion fuel, the neutrons come out in a preferred
direction --- and we built the first tool inside a standard reactor-physics code to simulate
that, then figured out when it actually helps.}
\vfill
Three worlds meet here:
\begin{itemize}
\item \key{Physics}: why aligned spins make neutrons directional.
\item \key{Engineering}: why a reactor designer cares where neutrons land.
\item \key{Computation}: how you simulate it honestly, and where cheap shortcuts lie to you.
\end{itemize}
\end{frame}

% ============================================================
\section{Part 1: Physics you need first}
\begin{frame}{Part 1}\Large Physics refreshers: fusion, spin, and magnetic confinement.\end{frame}

\begin{frame}{Fusion, and the neutron we care about}
The reaction in a D--T reactor:
\[ \underbrace{{}^2\mathrm{D}}_{\text{deuteron}} + \underbrace{{}^3\mathrm{T}}_{\text{triton}}
\;\rightarrow\; \underbrace{{}^4\mathrm{He}}_{3.5\,\text{MeV}} +
\underbrace{n}_{\key{14.1\,\text{MeV}}} \]
\begin{itemize}
\item The \key{neutron} carries 14.1 MeV (\,$1\,\text{MeV}=1.6\times10^{-13}$ J\,). It is
\key{uncharged}, so magnetic fields do \emph{not} bend it --- once born it flies straight
until it hits something.
\item That neutron is the entire product we track: it breeds tritium, heats the blanket, and
damages the walls. Everything downstream is ``where do the neutrons go?''
\end{itemize}
\end{frame}

\begin{frame}{Spin, in one slide (quantum refresher)}
\key{Spin} is an intrinsic angular momentum of a particle --- think of it as a tiny built-in
arrow that, in a magnetic field, can only point in a few discrete directions.
\begin{itemize}
\item The \key{deuteron} has spin 1: its projection on $\bhat$ can be $+1,\,0,\,-1$
(call the fractions $d_+,d_0,d_-$).
\item The \key{triton} has spin $\tfrac12$: projection $+\tfrac12$ or $-\tfrac12$
(fractions $t_+,t_-$).
\end{itemize}
Normally the fuel is \key{unpolarized}: all projections equally likely, arrows point every
which way. \key{Polarizing} the fuel means preparing the spins so specific projections
dominate. That is a knob a reactor designer could, in principle, turn.
\end{frame}

\begin{frame}{Magnetic confinement (E\&M refresher)}
A hot fusion plasma is charged (ions + electrons). From the Lorentz force
$\mathbf{F}=q\,\mathbf{v}\times\mathbf{B}$, a charged particle \key{spirals around magnetic
field lines} --- free to move along $\bhat$, tightly bound across it.
\begin{itemize}
\item So if you can make field lines that close on themselves inside a donut (torus), you
trap the plasma. That is a magnetic confinement device.
\item The field lines wrap around on nested \key{flux surfaces} --- like the layers of an
onion. Each surface has (roughly) constant temperature and density.
\item Key point for us: \key{at every point there is a well-defined field direction $\bhat$}.
The polarized neutrons are emitted relative to \emph{that local $\bhat$}.
\end{itemize}
\end{frame}

\begin{frame}{Two ways to build the donut}
\begin{columns}
\begin{column}{0.5\textwidth}
\key{Tokamak}
\begin{itemize}\footnotesize
\item Axisymmetric: looks the same at every toroidal angle $\phi$.
\item Needs a big electric current in the plasma to twist the field.
\item Simple shape; the current can cause disruptions.
\end{itemize}
\end{column}
\begin{column}{0.5\textwidth}
\key{Stellarator}
\begin{itemize}\footnotesize
\item Fully 3D: the cross-section twists and changes as you go around.
\item The twist is built into fancy 3D coils --- no plasma current needed, so it is steadier.
\item Much harder to design and build; the geometry is genuinely three-dimensional.
\end{itemize}
\end{column}
\end{columns}
\vfill
\key{Axisymmetry} (``same at every $\phi$'') is the single most important distinction for us:
it is what makes a tokamak easy to compute and a stellarator hard.
\end{frame}

% ============================================================
\section{Part 2: Spin-polarized fusion}
\begin{frame}{Part 2}\Large Why aligned spins make neutrons come out sideways.\end{frame}

\begin{frame}{The cross section: how likely, and in what direction}
A \key{cross section} $\sigma$ measures how likely a reaction is (units of area --- a bigger
``target''). The \key{differential cross section} $\dfrac{d\sigma}{d\Omega}$ adds
\emph{direction}: how likely to emit into a small solid angle $d\Omega$ around some
direction.
\begin{itemize}
\item \key{Solid angle} $\Omega$ = the 3D version of an angle; the full sphere is $4\pi$.
\item For unpolarized fuel, emission is \key{isotropic}: same in every direction,
$\frac{d\sigma}{d\Omega}=\frac{\sigma}{4\pi}$.
\item For \emph{polarized} fuel, it is not --- and that is the whole point.
\end{itemize}
\end{frame}

\begin{frame}{The master formula (Schwartz 2025)}
With $\theta_B$ = angle between the emitted neutron and the local field $\bhat$:
\[ \frac{d\sigma}{d\Omega} = \frac{\sigma_0}{2\pi}\Big[\tfrac34\,a\,\sin^2\theta_B
+ \big(\tfrac23 b + \tfrac13 c\big)\big(\tfrac14+\tfrac34\cos^2\theta_B\big)\Big] \]
The three numbers $\key{(a,b,c)}$ are set by how you polarized the fuel, and $a+b+c=1$:
\[ a = d_+t_+ + d_-t_-,\qquad b = d_0,\qquad c = d_+t_- + d_-t_+ .\]
Unpolarized fuel is $(\tfrac13,\tfrac13,\tfrac13)$. Read the formula as: \key{a bit of
``$\sin^2$'' shape plus a bit of ``$\cos^2$'' shape}, weighted by how you prepared the spins.
\end{frame}

\begin{frame}{What $\sin^2\theta_B$ and $\cos^2\theta_B$ look like}
\begin{itemize}
\item $\sin^2\theta_B$: \key{zero along $\bhat$, maximum perpendicular} to it. Neutrons
squirt out sideways, in the plane perpendicular to the field. Call this the \key{A / ``perp''}
shape.
\item $\tfrac14+\tfrac34\cos^2\theta_B$: \key{maximum along $\bhat$}, minimum sideways.
Neutrons prefer to go along the field. Call this the \key{B/C / ``par''} shape.
\end{itemize}
\key{Legendre reminder:} $\cos^2\theta = \tfrac13(1+2P_2(\cos\theta))$ where
$P_2(x)=\tfrac12(3x^2-1)$ is the 2nd Legendre polynomial. So both shapes are ``isotropic plus
a $P_2$ wobble.''
\end{frame}

\begin{frame}{Two knobs, not one --- keep them separate}
The formula secretly contains \key{two independent things}:
\begin{enumerate}
\item \key{Direction (shape).} Which way the neutrons prefer to go. Controlled by one number
\[ w(\theta_B)\propto 1 + a_2\,P_2(\cos\theta_B),\qquad a_2\in[-1,+1].\]
$a_2=-1$: pure perp. $a_2=+1$: pure par. $a_2=0$: isotropic.
\item \key{Rate (how many).} The total number of neutrons, via
$\eta = a + \tfrac23 b + \tfrac13 c$.
\end{enumerate}
\key{Design lesson:} in code, sample the \emph{direction} from the pdf and multiply the
\emph{count} by $\eta$ separately. Mixing them is the classic bug.
\end{frame}

\begin{frame}{The three modes, in plain words}
\begin{tabular}{lll}
\key{Mode} & \key{Rate $\eta$} & \key{What it does}\\\hline
Unpolarized & $2/3$ & baseline, neutrons everywhere\\
A (perp)    & $1$ (\,$+50\%$\,) & most power; neutrons sideways to $\bhat$\\
B (par)     & $2/3$ (same) & same power; neutrons along $\bhat$\\
C (par)     & $1/3$ (\,$-50\%$\,) & half the power; same direction as B\\
\end{tabular}
\vfill
\key{Surprise (and a correction to common lore):} B and C point the \emph{same} way ---
they differ only in \emph{how many}. C is \emph{not} isotropic. The $+50\%$ reactivity boost
(Kulsrud 1982) lives in mode A.
\end{frame}

% ============================================================
\section{Part 3: The geometry problem}
\begin{frame}{Part 3}\Large Flux surfaces, VMEC, and why stellarators are hard.\end{frame}

\begin{frame}{Describing the plasma shape: flux coordinates}
Instead of $(x,y,z)$ we use \key{flux coordinates} $(\rho,\theta,\zeta)$:
\begin{itemize}
\item $\rho$: which onion layer (0 = center, 1 = edge / \key{last closed flux surface}).
\item $\theta$: the short way around the donut (poloidal angle).
\item $\zeta$: the long way around the donut (toroidal angle).
\end{itemize}
A computer code called \key{VMEC} solves the plasma equilibrium and hands you, for every
$(\rho,\theta,\zeta)$: the real position $(R,Z,\phi)$, the field direction $\bhat$, and the
\key{Jacobian} $\sqrt{g}$.
\end{frame}

\begin{frame}{What $\sqrt{g}$ is, and why it matters (multivar calc)}
When you change variables in an integral you multiply by the Jacobian determinant:
\[ \int f\,dx\,dy\,dz = \int f\,\underbrace{\sqrt{g}}_{\text{Jacobian}}\,d\rho\,d\theta\,d\zeta.\]
\begin{itemize}
\item $\sqrt{g}$ tells you how much \emph{real volume} sits in each little
$(\rho,\theta,\zeta)$ cell. On a distorted 3D flux surface some cells are fat, some thin.
\item To place fusion births correctly you must weight by \key{reactivity $\times \sqrt{g}$}
--- otherwise you put too many neutrons where the coordinate grid is dense but the real
volume is small. Getting this right is a big part of ``rigorous.''
\end{itemize}
\end{frame}

\begin{frame}{Quasi-symmetry: QA and QH}
Stellarators are 3D, but a clever subclass \key{hides a hidden symmetry in the field
strength} $|\mathbf{B}|$ (not in the shape). Particles then confine almost as well as in a
tokamak.
\begin{itemize}
\item \key{QA} (quasi-axisymmetric): $|\mathbf{B}|$ looks tokamak-like.
\item \key{QH} (quasi-helical): $|\mathbf{B}|$ has a helical symmetry.
\item \key{nfp} = number of field periods (the shape repeats nfp times around).
\item \key{$\iota$} (iota) = rotational transform = how much a field line twists poloidally
per toroidal lap.
\end{itemize}
Database used: \key{QUASR}, a public catalogue of $\sim\!370{,}000$ optimized stellarators.
\end{frame}

\begin{frame}{The killer subtlety: symmetry $\neq$ gentle shape}
\qq{Good quasi-symmetry is \emph{bought} with strong 3D boundary shaping.}
So the devices with the best physics ($|\mathbf{B}|$ symmetry) are often the most
geometrically distorted. This anti-correlation is why ``the plasma looks almost
axisymmetric'' is a bad predictor of anything --- and it foreshadows why cheap analytic
methods struggle exactly on the good devices.
\vfill
We measure geometric distortion with \key{$\varepsilon_{\rm eff}$}:
\[ \varepsilon_{\rm eff} = \frac{\text{helical excursion of the source}}{\text{source-to-wall standoff}}.\]
Small = gentle, large = strongly shaped. Real optimized QAs sit at
$\varepsilon_{\rm eff}\approx 1$--$2.6$.
\end{frame}

% ============================================================
\section{Part 4: The engineering problem}
\begin{frame}{Part 4}\Large Neutron wall loading, breeding, and why anyone cares.\end{frame}

\begin{frame}{Where the neutrons land: the first wall}
\begin{itemize}
\item The \key{first wall} is the innermost solid surface facing the plasma. Behind it sits
the \key{blanket}.
\item \key{Neutron wall loading (NWL)} = neutron power hitting the wall per unit area
(MW/m$^2$). Too much $\Rightarrow$ the wall erodes, the magnets behind it get damaged.
\item \key{Peaking factor} $\mathrm{PF} = \dfrac{\max \mathrm{NWL}}{\text{average }\mathrm{NWL}}$.
PF $=1$ is perfectly uniform; big PF means a nasty hot spot. Designers want PF low.
\item \key{Inboard} = the inner (small-radius, ``center-stack'') wall, usually the tightest,
most shielding-starved real estate. \key{Outboard} = outer wall, roomier.
\end{itemize}
\end{frame}

\begin{frame}{Breeding: the reactor must make its own fuel}
Tritium is rare, so the reactor breeds it in the blanket:
\[ n + {}^6\mathrm{Li} \rightarrow {}^4\mathrm{He} + \mathrm{T} + 4.8\,\text{MeV}.\]
\begin{itemize}
\item \key{TBR} (tritium breeding ratio) = tritium atoms bred per fusion neutron. You need
\key{TBR $> 1$} (about 1.15 with margin) or the reactor starves.
\item \key{FLiBe} = a molten salt (Li, Be, F) used as breeder+coolant; \key{Be} multiplies
neutrons via $(n,2n)$.
\item Our finding: TBR \emph{per neutron} barely changes with polarization, so steering the
neutrons around does \emph{not} break breeding self-sufficiency. Good news.
\end{itemize}
\end{frame}

\begin{frame}{The dream: steer neutrons with spin}
\qq{If polarization aims neutrons, maybe we can aim them \emph{away} from the fragile
inboard wall and lightly-shielded magnets --- buying component lifetime for free.}
This is the motivating hope, and it is exactly what SPF does in spherical tokamaks (Bae 2025:
$+68\%$ magnet lifetime). The scientific question of this project:
\key{does it also work for stellarators, and how much of the promised benefit is real once
you account for everything?}
\end{frame}

% ============================================================
\section{Part 5: Monte Carlo and the source}
\begin{frame}{Part 5}\Large Simulating neutrons honestly.\end{frame}

\begin{frame}{Monte Carlo transport in 90 seconds}
\key{Monte Carlo} = simulate millions of individual neutrons with random numbers, then
average. \key{OpenMC} is the open-source code we use.
\begin{itemize}
\item Each neutron is \key{born} (the ``source''), then flies, maybe \key{scatters},
maybe gets absorbed, until it dies. Tally what hits the wall.
\item \key{Free-streaming / first-flight}: pretend neutrons never scatter --- they fly
straight from birth to wall. A useful idealization that isolates the source geometry.
\item \key{Statistical error} shrinks like $1/\sqrt{N}$: to resolve a $1\%$ effect you need
$\sim\!10^4$ counts \emph{per bin}. This is why small effects are expensive in MC.
\end{itemize}
\end{frame}

\begin{frame}{The ``source'' is the thing we built}
The \key{source} decides where each neutron is born and which way it points. Ours must:
\begin{enumerate}
\item pick a birth position (weighted by reactivity $\times\sqrt{g}$);
\item look up the local field $\bhat$ there;
\item sample a direction from $w(\theta_B)\propto 1+a_2 P_2(\cos\theta_B)$ \emph{about that
$\bhat$};
\item rotate that direction into lab coordinates.
\end{enumerate}
Doing step 2--3 \key{on the fly, per neutron, in the local field frame} is what makes it
``native'' and general --- and is why it works for a twisting stellarator field where $\bhat$
changes everywhere.
\end{frame}

\begin{frame}{How you sample a direction (the computation)}
\begin{itemize}
\item \key{Inverse-CDF sampling}: to draw from a pdf $p(x)$, invert its cumulative
$u=\int_{-\infty}^x p$; feed a uniform random $u\in[0,1)$, get $x$. Exact and fast.
\item \key{Rejection sampling}: throw darts under an easy envelope, keep those under $p$.
Slower but foolproof --- we use it to \emph{cross-check} the inverse-CDF.
\item \key{Stratified sampling}: split the mixture (perp vs par) and sample each exactly ---
gives exact draws and makes per-mode testing trivial.
\item \key{Thread safety}: OpenMC runs many neutrons in parallel threads. If the source
stored mutable state it would corrupt. Rule: set everything at construction, then
\key{never write} --- ``\texttt{const}.''
\end{itemize}
\end{frame}

\begin{frame}{A verification trick worth stealing: linearity in $a_2$}
Because $P_2$ enters \emph{linearly}, the wall load is exactly linear in $a_2$:
\[ \mathrm{NWL}(a_2) = \mathrm{NWL}_{\text{unpol}} + a_2\big(\mathrm{NWL}_{\text{par}}-\mathrm{NWL}_{\text{unpol}}\big).\]
Set $a_2=+1$ and $a_2=-1$ and add:
\[ \boxed{\;\mathrm{NWL}_{\text{perp}} + \mathrm{NWL}_{\text{par}} = 2\,\mathrm{NWL}_{\text{unpol}}\;}\]
This must hold \key{bin by bin}, with \emph{no analytic reference needed}. If it fails, the
sampler is wrong. Cheap, exact, self-checking --- the kind of test reviewers trust.
\end{frame}

% ============================================================
\section{Part 6: What we actually found}
\begin{frame}{Part 6}\Large The results --- and the honest caveats.\end{frame}

\begin{frame}{Finding 1: where the cheap analytic method lies}
The competing cheap method treats each source filament and computes flux to a wall patch by a
\key{point-patch} formula. On a real QA it agrees with MC to $\le 1.8\%$ \emph{except} on the
sharply curved inboard wall ($4$--$6\%$ off).
\begin{itemize}
\item A controlled test shows the analytic error grows like \key{(shaping)$^2$} and is
$\sim\!3\times$ worse on the convex inboard wall.
\item At realistic shaping ($\varepsilon_{\rm eff}\approx 2.3$) the analytic \emph{quadrature
itself} breaks (34\% error, negative loads). \key{This regime is Monte-Carlo-only.}
\end{itemize}
That is the intellectual justification for using a transport code at all.
\end{frame}

\begin{frame}{Finding 2: scattering dilutes the steering $\sim 3\times$}
\qq{The single most important honesty result.}
\begin{itemize}
\item Free-streaming says polarization shifts the inboard load by $\pm 40\%$.
\item Put in a \emph{real blanket} with scattering on, and it is only $\pm 13$--$15\%$.
\item Why: neutrons that enter the wall \key{scatter back} into the cavity as a
direction-blind diffuse fog, washing out the aim.
\end{itemize}
Geometric and analytic steering numbers are \key{optimistic by $\sim\!3\times$}, and
\emph{only a full transport calc reveals it}. The analytic theory structurally cannot.
\end{frame}

\begin{frame}{Finding 3: rate vs steering --- ``B is the sweet spot''}
Restoring the true neutron count $\eta$ per mode (vs unpolarized $=1$):
\begin{center}\small
\begin{tabular}{lcccc}
& fusion rate & center-stack load & outboard & TBR/neutron\\\hline
A & 1.50 & \key{1.70} & 1.36 & 1.14\\
B & 1.00 & \key{0.85} & 1.09 & 1.16\\
C & 0.50 & 0.42 & 0.55 & 1.16\\
\end{tabular}
\end{center}
\begin{itemize}
\item A: more power, but aims \emph{into} the center stack $\Rightarrow$ load compounds to
$1.70\times$.
\item \key{B: relieves the center stack ($0.85\times$) at no power or breeding cost} --- the
free lunch.
\item C: relieves the most, but you paid half your power for it.
\end{itemize}
\end{frame}

\begin{frame}{Finding 4: strong local actuator, weak global flattener}
\qq{Polarization moves a lot of load \emph{locally} (40--70\% control at a point) but barely
lowers the overall \emph{peaking factor}.}
Why: the peaking is set by a narrow \key{near-field hot spot} (the wall's closest approach to
the plasma). Angular redistribution reshapes the broad load but \key{cannot move that spike}.
\begin{itemize}
\item Example device: optimal polarization drops PF only $1.78\to1.49$.
\item \key{Two independent methods agree} (our MC \emph{and} the analytic companion): SPF is a
scalpel, not a flattening iron.
\end{itemize}
This reframes what SPF is \emph{for}: local fluence tailoring on critical parts, not global
peaking reduction.
\end{frame}

% ============================================================
\section{Part 7: The analytic companion}
\begin{frame}{Part 7}\Large anarrima: doing it with calculus instead of dice.\end{frame}

\begin{frame}{Why an analytic method at all}
Monte Carlo is noisy and slow for small effects and gives no exact derivatives. An analytic
formula is instant, noise-free, and \key{differentiable} (you can optimize with it). The
catch: it only works where the geometry is gentle. anarrima is that analytic method for
polarized NWL.
\end{frame}

\begin{frame}{Near-singular integrals and ``singularity subtraction''}
The wall-load integral looks like $\displaystyle\int \frac{f(\phi)}{D(\phi)^{3/2}}\,d\phi$,
where $D$ = squared distance from source to wall. When the wall nearly grazes the source,
$D\to 0$ and the integrand spikes --- ordinary numerical integration
(\key{Gauss--Legendre quadrature}) needs thousands of points and still fails.
\begin{itemize}
\item \key{Trick}: because $D(\phi)$ is a trig polynomial, you can find the spike's location,
depth, and curvature \key{in closed form}, subtract an exactly-integrable model of it, and
integrate the smooth leftover with \key{12 points} instead of 6000.
\item Kept \key{differentiable}, so you still get exact gradients. This is the freshest
numerical-methods contribution.
\end{itemize}
\end{frame}

\begin{frame}{Autodiff, elliptic integrals, and the two thresholds}
\begin{itemize}
\item \key{Automatic differentiation (autodiff / JAX)}: the computer differentiates your code
exactly, for free --- gold for optimization. anarrima is built to be autodiff-clean.
\item \key{Carlson elliptic integrals}: the special functions the toroidal geometry produces;
we implement them with exact hand-coded derivatives.
\item \key{Two $\varepsilon_{\rm eff}$ thresholds}: the \emph{series} converges only for
$\varepsilon_{\rm eff}\lesssim \tfrac12$; the \emph{quadrature} survives to $\sim\!2$. They
are different numbers --- a genuine, citable clarification.
\item \key{gc$_2$ correction}: to compute emission about the \emph{unit} field $\bhat=\mathbf
B/|\mathbf B|$, you must differentiate the normalization $|\mathbf B|$ too; the original
derivation dropped a 2nd-order piece. Restoring it pushes the polarized series from 2nd to
3rd order accuracy. (That was your ``gc2 improvement.'')
\end{itemize}
\end{frame}

\begin{frame}{A clean side result: the wall is a low-pass filter}
Question: does tokamak \key{field ripple} (the field wobble from having a finite number
$N$ of coils) show up on the wall?
\[ \text{transmission of harmonic }N \;\sim\; e^{-N\,a/R_0}. \]
\begin{itemize}
\item The line-of-sight integration \key{blurs out} high-frequency source structure. For a
real tokamak $N a/R_0\approx 6$, so the ripple is suppressed $\sim\!100\times$ --- the wall
signal is $\sim\!50$ parts per million. Negligible.
\item \key{Bonus insight}: this same blurring is \emph{why} ``plasma shape'' metrics failed
to predict wall behavior --- the wall literally cannot see high-harmonic shape. One mechanism
explains several dead ends.
\end{itemize}
\end{frame}

% ============================================================
\section{Part 8: Predictors and statistics}
\begin{frame}{Part 8}\Large What we tried to predict, and what honesty demands.\end{frame}

\begin{frame}{The hunt for a cheap predictor}
Dream: a single number from the plasma shape that predicts SPF benefit, so you skip the
expensive transport. Attempts:
\begin{itemize}
\item \key{Nematic order parameter} $Q$ (borrowed from liquid-crystal physics: how aligned a
set of direction vectors are) --- \key{did not predict}.
\item A ``C-term'' shape metric --- wrong order, \key{did not predict}.
\item \key{Field-direction coherence} (mean resultant length of $\bhat$ in the right frame):
sharper --- it cleanly separates QA from QH, but the scalar alone is weak.
\end{itemize}
Best current predictor is unglamorous: \key{available headroom} --- devices with a big hot
spot have more to gain. And the low-pass slide explains \emph{why} shape metrics can't work.
\end{frame}

\begin{frame}{Statistics honesty (why $n$ matters)}
\begin{itemize}
\item \key{Spearman correlation} $\rho$: rank-based measure of ``do these two move together.''
\item With only \key{$n=5$} devices, $\rho=\pm 0.6$ is \emph{within noise} --- you cannot
conclude anything. We said so, and launched a bigger batch ($n\approx 25$) to actually decide.
\item \key{Statistical power}: the ability to detect a real effect. Small $n$ = low power =
easy to fool yourself. Reporting the null honestly (``we could not tell'') is a result, not a
failure.
\end{itemize}
This honesty --- reporting the $\sim\!3\times$ dilution, the failed predictors, the
$n$-too-small caveats --- is what makes the work trustworthy.
\end{frame}

% ============================================================
\section{Part 9: The paper}
\begin{frame}{Part 9}\Large How it all becomes one story.\end{frame}

\begin{frame}{What is genuinely new (and what is not)}
\begin{itemize}
\item \key{New}: first \emph{native, on-the-fly, general} polarized fusion source in a
production MC code; first \emph{polarized stellarator} source; first stellarator source
\emph{in OpenMC}.
\item \key{Not first}: ``a polarized source in OpenMC'' (Bae 2025, but precomputed, one
tokamak); ``a VMEC stellarator source'' (Lyytinen 2024, but Serpent2, isotropic, geometry-
focused). We cite these up front and differentiate.
\item \key{Precedent}: Peterson's tokamak-source paper shows a native+validated tool is
publishable even as an increment --- and we clear that bar more easily (no prior stellarator
source in OpenMC; we add polarization).
\end{itemize}
\end{frame}

\begin{frame}{The paper's spine}
\qq{A native spin-polarized source for OpenMC --- and where/why Monte Carlo is necessary for
polarized stellarator wall loading.}
\begin{enumerate}\small
\item Physics + the native sampler.
\item StellaratorSource from a verified equilibrium.
\item Verification (incl.\ the linearity identity).
\item Where MC is necessary: analytic breakdown; the $3\times$ scattering dilution; rate-vs-
steering; the weak-global/strong-local finding.
\item Analytic companion (anarrima): singularity subtraction, differentiable quadrature,
$\varepsilon_{\rm eff}$ thresholds --- and it independently confirms the main finding.
\end{enumerate}
Genre: a \key{methods + honest-findings} paper. It does not need a blockbuster discovery; it
needs a correct, validated, useful capability, characterized honestly.
\end{frame}

\begin{frame}{The five things to remember}
\begin{enumerate}
\item Aligned spins $\Rightarrow$ directional neutrons: one knob for \key{direction} ($a_2$),
one for \key{rate} ($\eta$).
\item You need the \key{real local field $\bhat$} from a \key{real equilibrium} --- that is
what makes the stellarator source rigorous and the polarization well-posed.
\item \key{Scattering dilutes steering $\sim\!3\times$}; only transport shows it.
\item SPF is a \key{local scalpel, not a global flattener} --- two methods agree.
\item Cheap analytic methods \key{break where the geometry is strong}; that boundary is the
paper's thesis.
\end{enumerate}
\end{frame}

\begin{frame}{}
\begin{center}\Large Enjoy the train ride.\\[1em]
\normalsize Questions this deck should now let you answer:\\
\emph{Why do polarized neutrons care about $\bhat$? Why is a stellarator harder than a
tokamak? Why is free-streaming optimistic? What is a peaking factor, and why can't
polarization fix it? What does ``singularity subtraction'' buy you?}
\end{center}
\end{frame}

\end{document}
